% Plantilla común para figuras redibujadas (fig-NN-desc.tex).
% Copiar entera y sustituir solo el cuerpo entre \begin{document} y \end{document}.
% Compila tal cual en Overleaf o con scripts/tikz2svg.sh.
%
% NO añadir la opción "tikz" a standalone: parte cada tikzpicture (y cada
% \chemfig) en una página distinta y el SVG solo tendría la primera.
\documentclass[border=4pt]{standalone}
\usepackage[T1]{fontenc}
\usepackage{amsmath,amssymb}
\usepackage[version=4]{mhchem}  % \ce{A + B -> C}
\usepackage{chemfig}            % estructuras y esquemas de reacción
\usepackage{tikz}
\usetikzlibrary{arrows.meta,positioning,calc,shapes.geometric,decorations.pathmorphing,decorations.markings}
\usepackage{pgfplots}           % gráficas (cinética, curvas de valoración...)
\pgfplotsset{compat=1.18}

\begin{document}
% --- cuerpo de la figura ---
\begin{tikzpicture}
  \draw[-Stealth] (0,0) -- (2,0) node[right] {$x$};
\end{tikzpicture}
% ---------------------------
\end{document}
